\documentclass[11pt,letterpaper]{article}
\usepackage[utf8]{inputenc}
\usepackage{microtype}
\usepackage{lmodern}
\usepackage[T1]{fontenc}
\usepackage{geometry}
\geometry{left=2cm,right=2cm,top=2cm,bottom=2cm}
\usepackage{parskip}
\setlength{\parindent}{0pt}
\usepackage{enumitem}
\setlist{nosep}
\usepackage{titlesec}
\titleformat{\section}{\normalfont\large\bfseries}{\thesection.}{0.5em}{}
\titleformat{\subsection}{\normalfont\normalsize\bfseries}{\thesubsection.}{0.5em}{}
\renewcommand{\baselinestretch}{1.05}

\begin{document}

% Page 1: Executive-style summary
\begin{center}
  {\Large\bfseries Executive Summary — Dynamics and Control of an Octocopter with Flexible Structure}\\[6pt]
  {\small (Condensed from attached manuscript and prior analysis)}
\end{center}

\vspace{6pt}

\section*{Motivation and contribution}
- The manuscript departs from the conventional rigid-body multirotor
paradigm by embedding geometric and material compliance into the
vehicle frame. Alternating bending and torsional viscoelastic joints
along each arm, plus inward-tilted rotor thrust vectors, aim to
improve lateral disturbance rejection and provide richer actuation
without adding actuators.
\\
- Main technical contribution: a first-principles, scalable
Euler--Lagrange modeling framework for an octocopter with four
flexible arms. The model assembles Jacobians, mass matrix,
Coriolis/centrifugal terms, gravity, elastic/damping forces, and
generalized rotor inputs for a system with 6 rigid-body DOF + 4n
structural DOF + rotor DOF.

\section*{Modeling and methods}
- Mechanical concept: each arm is modeled as a chain of $n$ single-DoF
viscoelastic joints (odd joints: torsion; even joints: bending); each
arm ends with two rotors whose thrusts are biased toward the center of
mass.
\\
- Kinematics and dynamics: detailed derivation of velocity and
rotation Jacobians, Euler--Lagrange derivation yielding
$M(q)\ddot q + V(q,\dot q) + G(q) = Q_{\mathrm{gen}}(q) -
Q_{\mathrm{damp}}(\dot q)$.  Elastic potential and viscous damping are
included per joint; rotor thrusts and torques map into generalized
forces via link Jacobians.
\\
- Implementation: modular MATLAB implementation able to scale with
$n$. A baseline PID controller acts on rigid-body base states;
flexible states are included as aggregate feedback terms.

\section*{Key results}
- The derivation yields modular closed-form expressions for Jacobians, the mass matrix, and generalized forces that scale with the number of links per arm.\\
- In ideal hover the baseline PID stabilizes the vehicle with small altitude deviation.\\
- Under a simulated sustained 5\,N horizontal disturbance (wind-like), the vehicle shows significant lateral and altitude deviations, indicating limited disturbance rejection for the baseline controller.\\
- Authors position the modeling framework as a foundation for future
control work (adaptive/robust controllers) and experimental
validation.

\section*{Limitations noted}
- Results are simulation-only; no hardware validation reported.\\
- Control design limited to rigid-body PID with aggregate flexible feedback; model-based or robust/adaptive controllers are not evaluated.\\
- Disturbance scenarios, sensitivity analyses, and quantitative
performance metrics are limited in scope.

\vfill
\hrule
\bigskip
\noindent (End of page 1)

\newpage

% Page 2: Prioritized list of improvements
\begin{center}
  {\Large\bfseries Prioritized List of Improvements}
\end{center}

\vspace{6pt}

\section*{Top priorities (must-address)}
\begin{enumerate}[left=0pt,label=\textbf{\arabic*.},wide=0pt]
  \item \textbf{Add experimental validation}
    \begin{itemize}
      \item Build a subscale prototype or bench-test single-arm assemblies: static stiffness, modal frequencies, and damping identification to validate viscoelastic joint parameters and Jacobian mappings.
      \item Perform simple flight tests (tethered/indoor motion-capture) for hover and gust response comparisons with the model.
    \end{itemize}

  \item \textbf{Implement model-based and robust/adaptive controllers}
    \begin{itemize}
      \item Implement at least one model-based controller (e.g., inverse dynamics / feedback linearization, nonlinear MPC) using the derived model, and an adaptive/robust alternative to handle parameter uncertainty.
      \item Assess active vibration suppression strategies (collocated damping, modal LQR) for flexible modes.
    \end{itemize}

  \item \textbf{Quantify performance with standardized metrics}
    \begin{itemize}
      \item Report numeric metrics for each scenario: RMS tracking error, peak deviation, settling time, actuator effort (integrated thrust/torque), and disturbance-rejection gain.
      \item Summarize comparisons in concise tables.
    \end{itemize}
\end{enumerate}

\section*{Important analysis and simulation extensions}
\begin{enumerate}[resume,left=0pt,label=\textbf{\arabic*.},wide=0pt]
  \item \textbf{Sensitivity and robustness analysis}
    \begin{itemize}
      \item Parameter sweeps over: number of links $n$, link stiffness/damping, rotor tilt angles, and mass/inertia variations.
      \item Monte Carlo simulations to quantify performance variability and robustness margins.
    \end{itemize}

  \item \textbf{Broader disturbance testing}
    \begin{itemize}
      \item Add stochastic gusts, impulsive gusts, multi-directional loads and frequency-varying disturbances. Report frequency-dependent disturbance rejection (Bode-like or disturbance-to-output gains).
    \end{itemize}

  \item \textbf{Reduced-order modeling and modal control}
    \begin{itemize}
      \item Derive modal truncation (dominant flexible modes) to enable lower-order controllers and analytical insights.
      \item Propose hierarchical control: rigid-body outer loop + modal inner loop for flexible dynamics.
    \end{itemize}
\end{enumerate}

\section*{Formal analysis and comparisons}
\begin{enumerate}[resume,left=0pt,label=\textbf{\arabic*.},wide=0pt]
  \item \textbf{Closed-loop stability and robustness guarantees}
    \begin{itemize}
      \item Provide local stability analysis (linearization at hover, closed-loop eigenvalues) and, where feasible, Lyapunov or input-to-state stability statements for controllers designed with the model.
    \end{itemize}

  \item \textbf{Benchmark against related designs}
    \begin{itemize}
      \item Quantitatively compare with tilted-rotor, fully-actuated, or seesaw-based designs under identical simulation scenarios to highlight strengths and trade-offs.
    \end{itemize}
\end{enumerate}

\section*{Presentation, reproducibility, and writing}
\begin{enumerate}[resume,left=0pt,label=\textbf{\arabic*.},wide=0pt]
  \item \textbf{Improve clarity and documentation}
    \begin{itemize}
      \item Move lengthy algebraic expansions to appendices or supplementary material; keep the main text focused on interpretation and key expressions.
      \item Add a nomenclature table, define the configuration vector $q$ early, and state units for all parameters.
      \item Make figures self-contained with labeled joint indices, rotor offsets, and coordinate frames.
    \end{itemize}

  \item \textbf{Share code and parameters}
    \begin{itemize}
      \item Provide the simulation code, parameter files, and example scripts (if permitted) or pseudocode to increase reproducibility. Include a parameter table (link lengths, stiffnesses, damping coefficients, masses, PID gains) and brief rationale for selections.
    \end{itemize}
\end{enumerate}

\bigskip
\noindent \textbf{Suggested incremental experiments and workflow}
\begin{itemize}
  \item Bench tests: measure single-arm static stiffness and first 2--3 modal frequencies; update model stiffness/damping coefficients.
  \item Controller development: start with linearization about hover, design LQR/H$_\infty$ controllers or modal LQR for prototyping; move to nonlinear MPC or adaptive schemes for improved performance.
  \item Reporting: add Monte Carlo uncertainty bands in plots and condensed tables summarizing controller comparisons.
\end{itemize}

\vfill
\hrule
\smallskip
\noindent \textit{Note: When preparing hardware tests or releasing code, ensure compliance with institutional and confidentiality policies.}
\end{document}
